\subsection{Число е как предел последовательности.}

Вспомним неравенство среднего геометрического и среднего арифметического.
\[\forall k: a_k>0\;\;\;\ \sqrt[n]{a_1\cdot a_2\cdot ... \cdot a_n} \leq \cfrac{a_1+a_2+...+a_3}{n} \]
Пусть $a_1 = a > 0,\;a_2=a_3=...=a_{n+1}=d>0$\\
$\sqrt[n+1]{a\cdot b^n} \leq \cfrac{a+nb}{n+1}$\\
$(1+\cfrac{1}{n})^k$ - сходится
\begin{tcolorbox}[title=Доказательства монотонности]
    \begin{enumerate}
        \item $(1+\cfrac{1}{n})^k \uparrow$\\
        Пусть $a = 1,\;b=1+\cfrac{1}{n}$\\
        $\sqrt[n+1]{a\cdot(1+\cfrac{1}{n})^n}\; < \cfrac{1+n(1+\cfrac{1}{n})}{n+1} = \cfrac{n+2}{n+1} = 1+\cfrac{1}{n+1} \uparrow$\\
	Возведём в степень $n+1$
        \begin{tcolorbox}
            \[ (1+\cfrac{1}{n})^n < (1+\cfrac{1}{n+1})^{n+1} \]
        \end{tcolorbox}
        \item $(1-\cfrac{1}{n})^k \uparrow$\\
        Пусть $a = 1,\;b=1-\cfrac{1}{n}$\\
        $\sqrt[n+1]{a\cdot(1-\cfrac{1}{n})^n}\; < \cfrac{1+n(1-\cfrac{1}{n})}{n+1} = \cfrac{n+1-1}{n+1} = 1-\cfrac{1}{n+1} \uparrow$\\
	Возведём в степень $n+1$
        \begin{tcolorbox}
            \[ (1-\cfrac{1}{n})^n < (1-\cfrac{1}{n+1})^{n+1} \]    
        \end{tcolorbox}
        \item $(1+\cfrac{1}{n})^{n+1} \downarrow$\\
		$(1+\cfrac{1}{n})^{n+1} = (\cfrac{1}{\frac{n}{n+1}})^{n+1} = \cfrac{1}{(1-\frac{1}{n+1})^{n+1} \uparrow } $
    \end{enumerate}
\end{tcolorbox}
$\forall k,m \;\;\; n=max\{k,m\}$\\
\[ (1+\cfrac{1}{k})^k \leq (1+\cfrac{1}{n})^n < (1+\cfrac{1}{n})^{n+1} \leq (1+\cfrac{1}{m})^{m+1} \Rightarrow \]
$(1+\cfrac{1}{n})^n$ - ограничена сверху $\Rightarrow$ сходится.

\begin{tcolorbox}
    \textbf{Определение} 
    \[e = \lim_{n\to \infty} (1+\cfrac{1}{n})^n \]
\end{tcolorbox}

\begin{tcolorbox}[title=Доказать/подумать]
    \begin{enumerate}
        \item \[ e = \lim_{n\to\infty} (1+\cfrac{1}{1!} +\cfrac{1}{2!} + ... + +\cfrac{1}{n!}) \]
	\item $(1+\cfrac{1}{n})^{n+p}$ Когда возрастает, когда убывает и что происходит при $p=\frac{1}{2}$
        \item Сколько слагаемых нужно взять, чтобы получить e с точностью $10^{-3}$\\
        \item Какое нужно взять n, чтобы получить e с точностью $10^{-3}$
    \end{enumerate}
\end{tcolorbox}

\textbf{Утверждение.} $e = \lim_{n\to\infty} \left(1 + \frac{1}{1!} + \frac{1}{2!} + \dots + \frac{1}{n!}\right)$

\begin{tcolorbox}[title=Доказательство]
	Пусть $x_n = \left(1 + \frac{1}{n}\right)^n$ и $y_n = \sum_{k=0}^n \frac{1}{k!}$. Известно, что $x_n \to e$. Докажем, что $y_n \to e$.
	\begin{enumerate}
		\item Разложим $x_n$ по биному Ньютона:
			$$x_n = 1 + n\cdot\frac{1}{n} + \frac{n(n-1)}{2!}\frac{1}{n^2} + \dots + \frac{n(n-1)\dots(n-k+1)}{k! n^k} + \dots + \frac{1}{n^n}$$
			Преобразуем $k$-й член:
			$$\frac{1}{k!} \cdot \frac{n}{n} \cdot \frac{n-1}{n} \cdot \dots \cdot \frac{n-k+1}{n} = \frac{1}{k!} \left(1 - \frac{1}{n}\right)\dots\left(1 - \frac{k-1}{n}\right)$$
		\item $x_n < y_n$
			\begin{itemize}
				\item Заметим, что каждый множитель в скобках вида $\left(1 - \frac{m}{n}\right) < 1$.
				\item Значит, каждый член разложения $x_n$ строго меньше соответствующего члена $y_n$ ($\frac{1}{k!}$).
				\item $x_n < y_n \Rightarrow \lim_{n\to\infty} x_n \leq \lim_{n\to\infty} y_n \Rightarrow e \leq \lim y_n$ (так как $y_n \uparrow$, предел существует).
			\end{itemize}
		\item $y_n \leq e$
			\begin{itemize}
				\item Зафиксируем произвольное $k < n$. Рассмотрим сумму первых $k+1$ слагаемых в разложении $x_n$:
				$$x_n > 1 + 1 + \frac{1}{2!}(1 - \frac{1}{n}) + \dots + \frac{1}{k!}(1 - \frac{1}{n})\dots(1 - \frac{k-1}{n})$$
				\item Перейдем к пределу при $n \to \infty$ (держа $k$ фиксированным):
				$$e \geq 1 + 1 + \frac{1}{2!} + \dots + \frac{1}{k!} = y_k$$
				\item Так как $y_k \leq e$ для любого $k \Rightarrow \lim_{k\to\infty} y_k \leq e$.
			\end{itemize}
		\item \textbf{Вывод:}
			\begin{itemize}
				\item Из п.2 имеем $e \leq \lim y_n$.
				\item Из п.3 имеем $\lim y_n \leq e$.
				\item Следовательно, $\lim y_n = e$.
			\end{itemize}
	\end{enumerate}
\end{tcolorbox}

\textbf{Задача.} Исследовать на монотонность последовательность $x_n = \left(1+\frac{1}{n}\right)^{n+p}$ в зависимости от параметра $p$.

\begin{tcolorbox}[title=Анализ монотонности]
	Рассмотрим функцию $f(x) = \left(1+\frac{1}{x}\right)^{x+p}$ для $x > 0$.
	\begin{enumerate}
		\item \textbf{Логарифмирование:}
			$$ \ln f(x) = (x+p) \ln\left(1+\frac{1}{x}\right) $$
		\item \textbf{Производная:}
			$$ (\ln f(x))' = \ln\left(1+\frac{1}{x}\right) + (x+p) \cdot \frac{1}{1+\frac{1}{x}} \cdot \left(-\frac{1}{x^2}\right) $$
			$$ (\ln f(x))' = \ln\left(1+\frac{1}{x}\right) - \frac{x+p}{x(x+1)} $$
		\item \textbf{Асимптотическое разложение (ряд Тейлора):}
			При $x \to \infty$ (малом $\frac{1}{x}$) используем $\ln(1+t) = t - \frac{t^2}{2} + o(t^3)$.\\
			Разложим дробь:
			$$ \frac{x+p}{x^2+x} = \frac{x+p}{x^2(1+\frac{1}{x})} = \frac{1}{x}\left(1+\frac{p}{x}\right)\left(1-\frac{1}{x} + o\left(\frac{1}{x^2}\right)\right) $$
			$$ = \frac{1}{x}\left(1 - \frac{1}{x} + \frac{p}{x}\right) = \frac{1}{x} - \frac{1-p}{x^2} + o\left(\frac{1}{x^3}\right) $$
			Подставим в производную:
			$$ (\ln f(x))' \approx \left(\frac{1}{x} - \frac{1}{2x^2}\right) - \left(\frac{1}{x} - \frac{1-p}{x^2}\right) $$
			$$ (\ln f(x))' \approx \frac{1}{x^2} \left(-\frac{1}{2} + 1 - p\right) = \frac{1}{x^2}\left(\frac{1}{2} - p\right) $$
	\end{enumerate}
\end{tcolorbox}

\begin{tcolorbox}[title=Выводы о монотонности]
	Знак производной при больших $n$ определяется знаком скобки $\left(\frac{1}{2} - p\right)$:
	\begin{itemize}
		\item Если $\frac{1}{2} - p > 0 \Rightarrow \mathbf{p < \frac{1}{2}}$:\\
			Последовательность \textbf{возрастает} (начиная с некоторого $n$).\\
			\textit{Частный случай $p=0$: классическая $(1+1/n)^n \uparrow e$.}
		\item Если $\frac{1}{2} - p < 0 \Rightarrow \mathbf{p > \frac{1}{2}}$:\\
			Последовательность \textbf{убывает} (начиная с некоторого $n$).\\
			\textit{Частный случай $p=1$: классическая $(1+1/n)^{n+1} \downarrow e$.}
	\end{itemize}
\end{tcolorbox}

\begin{tcolorbox}[title=Случай p = 1/2]
	При $p=\frac{1}{2}$ член с $1/x^2$ обнуляется. Нужно рассмотреть следующий член разложения ($1/x^3$).
	\begin{itemize}
		\item $\ln(1+1/x) \approx \frac{1}{x} - \frac{1}{2x^2} + \frac{1}{3x^3}$
		\item Дробь $\frac{x+0.5}{x(x+1)}$ раскладывается как $\frac{1}{x} - \frac{0.5}{x^2} + \frac{0.5}{x^3}$
		\item Разность:
		$$ \left(\frac{1}{3} - \frac{1}{2}\right)\frac{1}{x^3} = -\frac{1}{6x^3} < 0 $$
	\end{itemize}
	\textbf{Вывод:} При $p=\frac{1}{2}$ последовательность \textbf{строго убывает} при всех $n$. Она сходится к числу $e$ быстрее, чем классические последовательности при $p=0$ или $p=1$.
\end{tcolorbox}

